% Appendix C. Uses the TikZ styles and the macro \SE defined at the top of Appendix B
% (ret, adv, thickline, thickadv, imp, dot, vert, lab).
\tikzset{
  blob/.style={draw,fill=gray!35,minimum width=0.6cm,minimum height=1.9cm,inner sep=0},
  cvert/.style={circle,draw,fill=white,inner sep=1.6pt} % charge (density) vertex
}
\newcommand{\cframe}[1]{% closed bubble of width #1
  \draw[thickline] (0,0)--(#1,0); \draw[thickadv] (0,-1.4)--(#1,-1.4);
  \draw[gray] (0,0)--(0,-1.4); \draw[gray] (#1,0)--(#1,-1.4);
  \node[vert] at (0,-0.7) {}; \node[vert] at (#1,-0.7) {};}
\newcommand{\crung}[1]{% one rung at x=#1
  \node[dot] at (#1,0) {}; \node[dot] at (#1,-1.4) {};
  \draw[imp] (#1,0)--(#1,-1.4);}

\newpage
\section*{Appendix C: Evaluation of the dressed bubble and the ladder}
\label{app:evaluation}

\noindent
In Appendix~B, we reduced the retarded--advanced part of the conductivity to three families of diagrams. Here we evaluate the first two: the dressed bubble (class~1, Sec.~C.1) and the ladder (class~2, Sec.~C.2). The maximally crossed series (class~3) is evaluated in the main text. We use only what B has set up: the dressed propagators and the pinched integral, Eq.~\eqref{eq:pinch}, and the expressions \eqref{eq:S0} and \eqref{eq:Sladder} for $S^{(0)}$ and $S^{(n)}_{\mathrm{ladder}}$.

We work at $E=E_{\mathrm F}$ with $k_{\mathrm F}\ell\gg1$, so that
$\gamma=1/2\tau\ll E_{\mathrm F}$. Momentum sums are done by the
rule
\[
\sum_{\mathbf p}F(\xi_{\mathbf p},\hat{\mathbf p})
\;\to\;
V\nu\int_{-\infty}^{\infty}\mathrm{d}\xi\;
\big\langle F(\xi,\hat{\mathbf p})\big\rangle_{\hat{\mathbf p}},
\]
where $\langle\cdots\rangle_{\hat{\mathbf p}}$ is the average over
the directions of $\mathbf p$ (over the solid angle $\Omega_d$),
the density of states \emph{per spin} $\nu$ and the velocity
$\mathbf v=v_{\mathrm F}\hat{\mathbf p}$ are evaluated at the Fermi
surface, and $d$ is the dimension. This neglects the variation of $\nu$ and $v$ over the
width $\gamma$ of the pinched region, which costs corrections of
relative order $1/E_{\mathrm F}\tau$. We also write
$D=v_{\mathrm F}^2\tau/d$; C.1 shows that it is the diffusion constant.

\subsection*{C.1\quad The dressed bubble: the Drude conductivity}

Class~1 is the bubble of B.5,
$S^{(0)}=\sum_{\mathbf k}v_{\mathbf kx}^2\,
\overline G^R(\mathbf k)\overline G^A(\mathbf k)$,
which has no rung and therefore no angular coupling between the two
lines: the direction of $\mathbf k$ enters only through $v_{\mathbf kx}^2$,
and the energy only through the pair $\overline G^R\overline G^A$.
So with the rule above the sum factorises,
\[
S^{(0)}
=V\nu\,\big\langle v_x^2\big\rangle_{\hat{\mathbf p}}
\int\mathrm{d}\xi\;\overline G^R\overline G^A
=V\nu\,\frac{v_{\mathrm F}^2}{d}\,\frac{\pi}{\gamma}
=V\nu\,\frac{v_{\mathrm F}^2}{d}\,2\pi\tau,
\]
where we used $\langle\hat p_x^2\rangle=1/d$ (the $d$ components
of $\hat{\mathbf p}$ are equivalent and $\sum_i\hat p_i^2=1$) and the
pinched integral $\pi/\gamma=2\pi\tau$, Eq.~\eqref{eq:pinch}. Inserting this in
the conductivity formula \eqref{eq:sigma},
\[
\sigma^{(0)}
=\frac{e^2g_s}{2\pi V}\,S^{(0)}
=e^2g_s\,\nu\,\frac{v_{\mathrm F}^2}{d}\,\tau
=e^2g_s\,\nu\,D .
\]
For free electrons $\nu(\epsilon)\propto\epsilon^{d/2-1}$, so the
density is $n=g_s\int_0^{E_{\mathrm F}}\nu\,\mathrm{d}\epsilon
=2g_s\nu E_{\mathrm F}/d$, and with
$v_{\mathrm F}^2=2E_{\mathrm F}/m$ we get
$g_s\nu v_{\mathrm F}^2/d=n/m$ in every dimension. Hence
\[
\sigma^{(0)}=\sigma_D=\frac{ne^2\tau}{m},
\]
the Drude conductivity. The form $e^2g_s\nu D$ above is the Einstein
relation. This is the dc value of the Drude form quoted in A.2, now with the
scattering time $\tau$ computed from $W_0$.

% ===== Figure C1: the bubble =========================================
\begin{figure}[ht]
\centering
\begin{tikzpicture}[scale=0.85]
 \cframe{4.4}
 \node[lab,left] at (-0.15,-0.7) {$v_x$};
 \node[lab,right] at (4.55,-0.7) {$v_x$};
 \node[lab,above] at (2.2,0.05) {$\overline G^R(\mathbf k)$};
 \node[lab,below] at (2.2,-1.45) {$\overline G^A(\mathbf k)$};
 \node at (5.7,-0.7) {$=$};
 \node[anchor=west] at (6.1,-0.7)
   {$V\nu\,\langle v_x^2\rangle_{\hat{\mathbf p}}\displaystyle\int\mathrm{d}\xi\;
   \overline G^R\overline G^A=V\nu\,\dfrac{v_{\mathrm F}^2}{d}\,2\pi\tau$};
\end{tikzpicture}
\caption{Class~1, the dressed bubble, and its evaluation: the sum over
$\mathbf k$ splits into an average over the direction of $\mathbf k$
(which gives $\langle v_x^2\rangle=v_{\mathrm F}^2/d$) and the
pinched energy integral, Eq.~\eqref{eq:pinch}.}
\label{fig:bubbleeval}
\end{figure}


\subsection*{C.2\quad Class 2: the ladder as a vertex correction}

Class~2 adds uncrossed rungs to the bubble. As in
Ref.~\cite{Coleman2015}, we sum them by absorbing the rungs into a
dressed current vertex, which obeys a Bethe--Salpeter equation, and we
solve that equation on the Fermi surface. The method does not need white noise, so we let the Gaussian pair
correlator depend on the momentum transfer,
$\overline{U_{\mathbf a\mathbf b}U_{\mathbf c\mathbf d}}
=(W(\mathbf a-\mathbf b)/V)\,\delta_{\mathbf a-\mathbf b,\,\mathbf d-\mathbf c}$,
rotation invariant, so that at the Fermi surface it depends only on the
angle $\theta$ between the two momenta; white noise is $W=W_0$. The
Born rate of B.3 becomes $1/\tau=2\pi\nu\langle W\rangle$, where
$\langle\cdots\rangle$ is the average over $\theta$ on the Fermi
surface.

We work in the uniform limit of A.1 ($\mathbf q\to0$ first, then
$\omega\to0$) and write a retarded--advanced pair as
$\mathcal P_{\mathbf p}=\overline G^R(\mathbf p,E_{\mathrm F}+\omega)\,
\overline G^A(\mathbf p,E_{\mathrm F})$. Between two rungs the two lines carry
the same momentum (B.4), so a ladder with $n$ rungs has $n+1$ pairs,
and the first terms are (Fig.~\ref{fig:ladderterms})
\[
\begin{aligned}
S^{(0)}&=\sum_{\mathbf p}v_x(\mathbf p)\,\mathcal P_{\mathbf p}\,v_x(\mathbf p),\\
S^{(1)}&=\frac1V\sum_{\mathbf p,\mathbf p'}
v_x(\mathbf p)\,\mathcal P_{\mathbf p}\,W(\mathbf p-\mathbf p')\,
\mathcal P_{\mathbf p'}\,v_x(\mathbf p'),\\
S^{(2)}&=\frac1{V^2}\sum_{\mathbf p,\mathbf p',\mathbf p''}
v_x(\mathbf p)\,\mathcal P_{\mathbf p}\,W(\mathbf p-\mathbf p')\,
\mathcal P_{\mathbf p'}\,W(\mathbf p'-\mathbf p'')\,
\mathcal P_{\mathbf p''}\,v_x(\mathbf p''),
\end{aligned}
\]
and $S^{(0)}$ is the bubble of C.1. Each further rung adds a momentum
sum, a weight $W/V$ and a pair.

% ===== Figure: the first terms =======================================
\begin{figure}[ht]
\centering
\begin{tikzpicture}[scale=0.55]
 \begin{scope}
  \cframe{3.6}
  \node[lab,above] at (1.8,0.05) {$\mathbf p$};
  \node[lab] at (1.8,-2.3) {$S^{(0)}$};
  \node at (4.6,-0.7) {$+$};
 \end{scope}
 \begin{scope}[xshift=5.6cm]
  \cframe{5.4}\crung{2.7}
  \node[lab,above] at (1.35,0.05) {$\mathbf p$};
  \node[lab,above] at (4.05,0.05) {$\mathbf p'$};
  \node[lab] at (2.7,-2.3) {$S^{(1)}$};
  \node at (6.4,-0.7) {$+$};
 \end{scope}
 \begin{scope}[xshift=13.2cm]
  \cframe{7.2}\crung{2.4}\crung{4.8}
  \node[lab,above] at (1.2,0.05) {$\mathbf p$};
  \node[lab,above] at (3.6,0.05) {$\mathbf p'$};
  \node[lab,above] at (6.0,0.05) {$\mathbf p''$};
  \node[lab] at (3.6,-2.3) {$S^{(2)}$};
  \node at (8.4,-0.7) {$+\ \cdots$};
 \end{scope}
\end{tikzpicture}
\caption{The bubble and the ladders with one and two uncrossed rungs.
Between rungs both lines carry the same momentum.}
\label{fig:ladderterms}
\end{figure}

We collect the bare current vertex and every sequence of rungs next to it
into $\Gamma_x(\mathbf p)$,
\begin{center}
\begin{tikzpicture}[scale=0.6]
 \node at (-0.6,-0.7) {$\Gamma_x=$};
 \begin{scope}[xshift=1.0cm]
  \draw[thickline] (0,0)--(2.7,0);
  \draw[thickadv] (0,-1.4)--(2.7,-1.4);
  \draw[gray] (2.7,0)--(2.7,-1.4);
  \node[cvert] at (2.7,-0.7) {};
  \node at (4.0,-0.7) {$+$};
 \end{scope}
 \begin{scope}[xshift=6.2cm]
  \draw[thickline] (0,0)--(4.6,0);
  \draw[thickadv] (0,-1.4)--(4.6,-1.4);
  \crung{1.0}
  \draw[gray] (4.6,0)--(4.6,-1.4);
  \node[cvert] at (4.6,-0.7) {};
  \node at (5.9,-0.7) {$+$};
 \end{scope}
 \begin{scope}[xshift=13.2cm]
  \draw[thickline] (0,0)--(6.6,0);
  \draw[thickadv] (0,-1.4)--(6.6,-1.4);
  \crung{1.0}\crung{3.0}
  \draw[gray] (6.6,0)--(6.6,-1.4);
  \node[cvert] at (6.6,-0.7) {};
  \node at (8.0,-0.7) {$+\ \cdots$};
 \end{scope}
\end{tikzpicture}
\end{center}
\[
\Gamma_x(\mathbf p)=v_x(\mathbf p)
+\frac1V\sum_{\mathbf p'}W(\mathbf p-\mathbf p')\,\mathcal P_{\mathbf p'}\,v_x(\mathbf p')
+\frac1{V^2}\sum_{\mathbf p',\mathbf p''}W(\mathbf p-\mathbf p')\,\mathcal P_{\mathbf p'}\,
W(\mathbf p'-\mathbf p'')\,\mathcal P_{\mathbf p''}\,v_x(\mathbf p'')+\cdots .
\]
Factoring the first rung out of every term after the bare vertex leaves
the same series again, so
\begin{equation}
\Gamma_x(\mathbf p)=v_x(\mathbf p)
+\frac1V\sum_{\mathbf p'}W(\mathbf p-\mathbf p')\,\mathcal P_{\mathbf p'}\,\Gamma_x(\mathbf p'),
\label{eq:vertex}
\end{equation}
the Bethe--Salpeter equation of the ladder (Fig.~\ref{fig:vertexeq}).
Closing the dressed vertex with the left bare vertex and the first pair
(Fig.~\ref{fig:vertexclosed}) gives the bubble plus all the ladders,
\[
S^{(0)}+\sum_{n\ge1}S^{(n)}_{\mathrm{ladder}}
=\sum_{\mathbf p}v_x(\mathbf p)\,\mathcal P_{\mathbf p}\,\Gamma_x(\mathbf p),
\]
so the class-2 contribution is
$\sum_{\mathbf p}v_x\mathcal P_{\mathbf p}\,[\Gamma_x-v_x]$.

% ===== Figure: the vertex equation ====================================
\begin{figure}[ht]
\centering
\begin{tikzpicture}[scale=0.75]
 \begin{scope}
  \draw[thickline] (0,0)--(3.1,0);
  \draw[thickadv] (0,-1.4)--(3.1,-1.4);
  \node[blob] at (3.4,-0.7) {$\Gamma_x$};
  \node[lab,above] at (1.5,0.05) {$\mathbf p$};
  \node at (4.5,-0.7) {$=$};
 \end{scope}
 \begin{scope}[xshift=5.3cm]
  \draw[thickline] (0,0)--(3.1,0);
  \draw[thickadv] (0,-1.4)--(3.1,-1.4);
  \draw[gray] (3.1,0)--(3.1,-1.4);
  \node[vert] at (3.1,-0.7) {};
  \node[lab,right] at (3.25,-0.7) {$v_x$};
  \node[lab,above] at (1.5,0.05) {$\mathbf p$};
  \node at (4.7,-0.7) {$+$};
 \end{scope}
 \begin{scope}[xshift=11.0cm]
  \draw[thickline] (0,0)--(6.1,0);
  \draw[thickadv] (0,-1.4)--(6.1,-1.4);
  \crung{1.0}
  \node[lab,above] at (0.45,0.05) {$\mathbf p$};
  \node[lab,above] at (3.3,0.05) {$\mathbf p'$};
  \node[lab,below] at (3.3,-1.45) {$\mathcal P_{\mathbf p'}$};
  \node[lab,left] at (0.9,-0.7) {$\dfrac{W}{V}$};
  \node[blob] at (6.4,-0.7) {$\Gamma_x$};
 \end{scope}
\end{tikzpicture}
\caption{The Bethe--Salpeter equation for the dressed current vertex
(shaded): the bare vertex, plus one rung, a pair at the new momentum
$\mathbf p'$, and again the dressed vertex.}
\label{fig:vertexeq}
\end{figure}

% ===== Figure: closing ===============================================
\begin{figure}[ht]
\centering
\begin{tikzpicture}[scale=0.8]
 \draw[thickline] (0,0)--(6.1,0);
 \draw[thickadv] (0,-1.4)--(6.1,-1.4);
 \draw[gray] (0,0)--(0,-1.4);
 \node[vert] at (0,-0.7) {};
 \node[lab,left] at (-0.15,-0.7) {$v_x$};
 \node[blob] at (6.4,-0.7) {$\Gamma_x$};
 \node[lab,above] at (3.0,0.05) {$\mathcal P_{\mathbf p}$};
\end{tikzpicture}
\caption{The bubble together with all its ladders:
$\sum_{\mathbf p}v_x\,\mathcal P_{\mathbf p}\,\Gamma_x(\mathbf p)$.}
\label{fig:vertexclosed}
\end{figure}


The pair is peaked at the Fermi surface. With $\mathbf q=0$ and
$\omega\tau\ll1$ its energy integral is the pinched value
$\int\mathrm{d}\xi\,\mathcal P=2\pi\tau/(1-\mathrm{i}\omega\tau)\simeq2\pi\tau$,
Eq.~\eqref{eq:pinch}. Writing $\frac1V\sum_{\mathbf p'}\to\nu\int\mathrm{d}\xi'\langle\cdots\rangle_{\hat{\mathbf p}'}$
and neglecting the energy dependence of $W$ and $v$, the vertex depends
only on the direction of $\mathbf p$ and
\[
\Gamma_x(\hat{\mathbf p})
=v_{\mathrm F}\hat p_x
+2\pi\nu\tau\,\big\langle W(\hat{\mathbf p}\!\cdot\!\hat{\mathbf p}')\,
\Gamma_x(\hat{\mathbf p}')\big\rangle_{\hat{\mathbf p}'} .
\]

The kernel $W(\hat{\mathbf p}\!\cdot\!\hat{\mathbf p}')$ is rotation
invariant, so a harmonic of degree $l$ is an eigenfunction (Funk--Hecke
theorem): $\langle W(\hat{\mathbf p}\!\cdot\!\hat{\mathbf p}')\,Y_l(\hat{\mathbf p}')\rangle
=w_l\,Y_l(\hat{\mathbf p})$ with $w_l=\langle W\,P_l(\cos\theta)\rangle$,
$P_1=\cos\theta$. The bare vertex $v_{\mathrm F}\hat p_x$ is a harmonic
with $l=1$, so $\Gamma_x=\lambda\,v_{\mathrm F}\hat p_x$ with
\[
\lambda=1+2\pi\nu\tau\,\langle W\cos\theta\rangle\,\lambda .
\]
Using $2\pi\nu\tau=1/\langle W\rangle$,
\begin{equation}
\lambda=\frac{1}{1-\langle W\cos\theta\rangle/\langle W\rangle}
=\frac{\langle W\rangle}{\langle W\,(1-\cos\theta)\rangle}
=\frac{\tau_{\mathrm{tr}}}{\tau},
\qquad
\frac1{\tau_{\mathrm{tr}}}=2\pi\nu\,\langle W\,(1-\cos\theta)\rangle .
\label{eq:lambda}
\end{equation}


Since $\Gamma_x=\lambda v_x$, the closed diagram of
Fig.~\ref{fig:vertexclosed} is $\lambda S^{(0)}$, the class-2 part is
$(\lambda-1)S^{(0)}$, and
\begin{equation}
\sigma=\lambda\,\sigma_D=\frac{ne^2\tau_{\mathrm{tr}}}{m}.
\label{eq:sigmatr}
\end{equation}
The ladder therefore replaces the scattering time by the transport
time: forward scattering ($\cos\theta\approx1$) relaxes the density but
not the current. These are the standard steps and result.

\paragraph{White noise.}
For $W=W_0$ we have $\langle W_0\cos\theta\rangle=0$, so $\lambda=1$,
$\tau_{\mathrm{tr}}=\tau$, and the class-2 contribution is zero: the
ladder adds nothing to $\sigma_D$ for isotropic white-noise disorder. We
see it term by term as well. The first rung produces
$\mathcal J\equiv\sum_{\mathbf p}v_x\mathcal P_{\mathbf p}=V\nu\,2\pi\tau\,\langle v_x\rangle=0$
at each end, because the pair does not depend on the direction of
$\mathbf p$ while $v_x$ is odd (Fig.~\ref{fig:onerung}); the longer
ladders contain the same factors. This is the result used in B.5 for the
white-noise model.

% ===== Figure: white noise =========================================
\begin{figure}[ht]
\centering
\begin{tikzpicture}[scale=0.85]
 \cframe{5.4}\crung{2.7}
 \draw[gray] (0.1,-1.95)--(0.1,-2.1)--(2.6,-2.1)--(2.6,-1.95);
 \draw[gray] (2.8,-1.95)--(2.8,-2.1)--(5.3,-2.1)--(5.3,-1.95);
 \node[lab] at (1.35,-2.55) {$\mathcal J=0$};
 \node[lab] at (4.05,-2.55) {$\mathcal J=0$};
\end{tikzpicture}
\caption{White noise: in the one-rung ladder each pair touches a current
vertex and gives $\mathcal J=\sum_{\mathbf p}v_x\mathcal P_{\mathbf p}=0$.}
\label{fig:onerung}
\end{figure}

Throughout we took $\mathbf q\to0$ first, at finite $\omega$. The
vanishing for white noise holds in this order; the opposite order
behaves differently, because the ladder sum has a pole at
$\mathbf q=\omega=0$ (next section).
\subsection*{C.3\quad The ladder as a diffuson}

The same ladder also dresses a vertex that carries no velocity: the
charge vertex associated with the density. Its sum is the diffuson.
This provides a useful completion of the class-2 calculation: it
shows how the ladder behaves at finite wave vector and frequency, and
why its diffusion pole does not produce a correction to the uniform
dc conductivity. We take white-noise disorder and retain the external
wave vector and frequency in the retarded--advanced pair,
\[
\mathcal P_{\mathbf p}(\mathbf q,\omega)
=\overline G^R(\mathbf p+\mathbf q,E_{\mathrm F}+\omega)\,
 \overline G^A(\mathbf p,E_{\mathrm F}),
\]
where the offset $\mathbf q$ is carried unchanged through each rung,
since the rung transfers equal and opposite momenta to the two
propagator lines (Eq.~\eqref{eq:pairrule}).

The bare charge vertex is $1$. For white-noise disorder the rung is
momentum independent, so the dressed charge vertex is also independent
of $\mathbf p$. The Bethe--Salpeter equation therefore reduces to
\[
\Gamma_c
=1+\frac{W_0}{V}\sum_{\mathbf p}
  \mathcal P_{\mathbf p}\Gamma_c
=1+W_0\Lambda\,\Gamma_c,
\qquad
\Lambda\equiv\frac1V\sum_{\mathbf p}\mathcal P_{\mathbf p}.
\]
Solving gives
\[
\Gamma_c=\frac{1}{1-W_0\Lambda}
=\sum_{m\ge0}(W_0\Lambda)^m.
\]
Thus each additional rung--pair unit contributes a factor
$W_0\Lambda$ to the geometric series.


For $q\ell\ll1$, we expand
$\xi_{\mathbf p+\mathbf q}\simeq
\xi_{\mathbf p}+\mathbf v\cdot\mathbf q$. Writing
$\xi=\xi_{\mathbf p}$, the pair becomes
\[
\mathcal P_{\mathbf p}
=\frac{1}{
(\omega-\mathbf v\cdot\mathbf q-\xi+\mathrm{i}\gamma)
(-\xi-\mathrm{i}\gamma)}.
\]
The poles in $\xi$ lie at
$\omega-\mathbf v\cdot\mathbf q+\mathrm{i}\gamma$
and $-\mathrm{i}\gamma$. Closing the contour in the upper half-plane,
as in Eq.~\eqref{eq:pinch}, gives
\[
\int\mathrm{d}\xi\,\mathcal P_{\mathbf p}
=\frac{2\pi\mathrm{i}}
{\omega-\mathbf v\cdot\mathbf q+2\mathrm{i}\gamma}
=\frac{2\pi\tau}
{1-\mathrm{i}\omega\tau+
 \mathrm{i}\tau\,\mathbf v\cdot\mathbf q}.
\]
Using the momentum-sum rule and expanding for
$|\omega|\tau\ll1$ and $q\ell\ll1$, the angular average gives
\[
\Lambda
=2\pi\nu\tau
\left\langle
\frac{1}{1-\mathrm{i}\omega\tau+
\mathrm{i}\tau\,\mathbf v\cdot\mathbf q}
\right\rangle_{\hat{\mathbf p}}
\simeq
2\pi\nu\tau
\left[1+\mathrm{i}\omega\tau-D\tau q^2\right].
\]
Here the term linear in $\mathbf v\cdot\mathbf q$ vanishes by angular
averaging, while
$\langle(\mathbf v\cdot\mathbf q)^2\rangle
=v_{\mathrm F}^2q^2/d=Dq^2/\tau$.
Using $W_0\,2\pi\nu\tau=1$ from Eq.~\eqref{eq:tau}, we obtain
\[
1-W_0\Lambda
\simeq\tau(Dq^2-\mathrm{i}\omega),
\]
and hence
\begin{equation}
\Gamma_c(\mathbf q,\omega)
=\frac{1}{1-W_0\Lambda}
=\frac{1}{\tau(Dq^2-\mathrm{i}\omega)}.
\label{eq:diffpole}
\end{equation}
This is the diffusion pole. It describes the long-wavelength,
low-frequency propagation of a density fluctuation, whose relaxation
rate is $Dq^2$. At $\mathbf q=\mathbf 0$ and $\omega=0$, each
rung--pair unit has weight $W_0\Lambda=1$, so the geometric series
becomes singular. The pole is therefore a consequence of the
conservation of particle number.

Why does this singularity not affect the conductivity of C.2? The
diffuson is a density mode, and a density mode is not excited by a
uniform field. A field with $\mathbf q=0$ drives a current but cannot
change the local density, since by the continuity equation
$\partial_t\rho_{\mathbf q}=-\mathrm{i}\,\mathbf q\!\cdot\!\mathbf j_{\mathbf q}$
vanishes at $\mathbf q=0$. The current vertex $v_x$ is odd in $\hat{\mathbf p}$, the pole
sits in the isotropic ($l=0$) channel, and the two meet only through
the factor $\mathcal J=\sum_{\mathbf p}v_x\mathcal P_{\mathbf p}$, which
is exactly zero at $\mathbf q=0$ and of first order in $\mathbf q$ at
small $\mathbf q$. The ladder with current vertices at both ends is
therefore proportional to
\[
\mathcal J^2\,\Gamma_c\;\propto\;\frac{q_x^2}{Dq^2-\mathrm{i}\omega},
\]
so the pole is always accompanied by a factor that vanishes with $q$. For
$\mathbf q\to0$ at fixed $\omega\neq0$ the denominator stays finite and the expression goes to zero, which is the order of limits shown A.1 for dc conductivity, and
the ladder leaves $\sigma_D$ unchanged apart from the
$\tau\to\tau_{\mathrm{tr}}$ replacement of C.2. What looked like an infinity multiplied by zero is
resolved by taking $\mathbf q\to0$ first, as the dc conductivity requires.
The opposite order ($\omega\to0$ at fixed $\mathbf q$) is a different
physical question: it asks how a static density modulation relaxes, and the answer, a
vanishing longitudinal current response, is the statement that a conserved density cannot
flow forever, not a property of the dc conductivity.

This pole starts to matter in two situations. The first is when the probe
itself carries a wave vector, so that $q$ is an external parameter
and $\Gamma_c$ is probed directly: the density response, the diffusion
of a charge packet, the longitudinal conductivity at finite $q$ and
$\omega$. There the pole appears physically, as the diffusive relaxation of density
at rate $Dq^2$, with $D$ the diffusion constant of C.1. The second, and for our paper
the decisive one, is when the pole momentum is not external but
internal, and is summed over. In the diffuson the pole carries the
momentum transferred by the probe, which we hold fixed and let go to
zero. In the maximally crossed series the pole carries $\mathbf Q=\mathbf k+\mathbf k'$,
the sum of the incoming and outgoing momenta; time-reversed paths pair
$\mathbf k'\simeq-\mathbf k$, so the current vertices at the two ends
are $v_x(\mathbf k)\,v_x(\mathbf k')\simeq-v_x^2(\mathbf k)$ and do
\emph{not} vanish as $\mathbf Q\to0$, and $\mathbf Q$ is an integration variable, not a probe.
No factor of $q^2$ cancels the pole. The correction to $\sigma$ is then
a sum over all $\mathbf Q$ of $-1/Q^2$-type terms, cut off at large $Q$ by
$1/\ell$, where the diffusion approximation fails, and at small $Q$ by the
sample size $L$ or by dephasing, which replaces $-\mathrm{i}\omega$ in the pole
by $1/\tau_\varphi$. Thus the same pole structure that is harmless for the
ladder produces the quantum correction of the crossed series. This is the
route to the scaling theory, and its dimension dependence follows from
how $\int\mathrm{d}^dQ/Q^2$ behaves at small $Q$: it converges in $d=3$, so the
correction is a small finite number set by the large-$Q$ cutoff; it diverges
logarithmically in $d=2$, so the correction grows without bound as the
infrared cutoff is removed; and it diverges as a power in $d=1$. The main
text carries this out for the crossed series, using the same three steps
as C.3.